

\chapter{Dictionary to the book: Chapter 1 notation}\label{app:dict:c1}
%=======================================================================

\renewcommand{\arraystretch}{1.2}
\begin{center}\footnotesize
\begin{tabular}{@{}ll@{\hspace{1.3em}}ll@{}}
\hline
Book & Here & Book & Here\\
\hline
(1.1)--(1.3) & Def.~\ref{def:lindep}, Thm.~\ref{thm:coord}
 & (1.79)--(1.84) & Thm.~\ref{thm:trace}\\
ip axioms & Def.~\ref{def:ip}
 & (1.85), (1.86) & Def.~\ref{def:cof}, Prop.~\ref{prop:I2cof}\\
(1.4)--(1.6) & Thm.~\ref{thm:cs}--Cor.~\ref{cor:tri}
 & (1.87)--(1.93) & Def.~\ref{def:tip}, Thm.~\ref{thm:tip}\\
(1.7)--(1.16) & Def.~\ref{def:onb}--Prop.~\ref{prop:comp}
 & (1.95) & Prop.~\ref{prop:posvec}\\
(1.17)--(1.19) & Thm.~\ref{thm:riesz}
 & \eqr{1.96}, (1.97) & Def.~\ref{def:diff}, Def.~\ref{def:landau}\\
(1.20)--(1.22) & Thm.~\ref{thm:cob}
 & (1.98), \eqr{1.99} & Thm.~\ref{thm:grad}\\
(1.23)--(1.30) & Def.~\ref{def:cross}--Prop.~\ref{prop:crosscomp}
 & \eqr{1.100}--(1.108) & Thm.~\ref{thm:dir}--Thm.~\ref{thm:cart}\\
(1.31)--(1.39) & Def.~\ref{def:det}--Cor.~\ref{cor:det6}
 & \eqr{1.109}--(1.115) & Thm.~\ref{thm:chain}--Ex.~\ref{ex:cartrecover}\\
(1.40)--(1.50) & Def.~\ref{def:tensor}--Ex.~\ref{ex:tp}
 & \eqr{1.116}--\eqr{1.125} & Ex.~\ref{ex:polar}, Ex.~\ref{ex:gradr}\\
(1.55)--(1.58) & Thm.~\ref{thm:tenstrans}
 & \eqr{1.126}--(1.134) & Def.~\ref{def:diffv}--Cor.~\ref{cor:dv}\\
(1.59)--(1.62) & Prop.~\ref{prop:transpose}, Prop.~\ref{prop:axial}
 & (1.135)--\eqr{1.145} & Ex.~\ref{ex:gu1}--Ex.~\ref{ex:gupolar}\\
(1.63)--(1.72) & Prop.~\ref{prop:prodinv}
 & (1.146)--(1.152) & Def.~\ref{def:divv}--Prop.~\ref{prop:divprod}\\
(1.73)--(1.77) & Prop.~\ref{prop:orth}
 & (1.153)--\eqr{1.162} & Ex.~\ref{ex:divpolar}\\
(1.78) & Def.~\ref{def:inv3}
 & (1.163)--(1.166) & Def.~\ref{def:curlv}--Rem.~\ref{rem:curlA}\\
sym/skw, sph/dev & Def.~\ref{def:splits}
 & \eqr{1.167}, \eqr{1.168} & Thm.~\ref{thm:div}, Thm.~\ref{thm:stokes}\\
\hline
\end{tabular}
\end{center}

\medskip\noindent
Lemma~\ref{lem:altform} (one dimensionality of alternating $3$ forms)
replaces the book's Appendix~A computations in
Propositions~\ref{prop:detforms},~\ref{prop:invwd} and
Theorem~\ref{thm:det}. All twenty eight chapter problems are solved in
\S\ref{sec:problems:c1}.



\chapter{Dictionary to the book: Chapter 2 notation}\label{app:dict:c2}
%=======================================================================

\renewcommand{\arraystretch}{1.2}
\begin{center}\scriptsize
\begin{tabular}{@{}ll@{\hspace{1.6em}}ll@{}}
\hline
Book & Here & Book & Here\\
\hline
\eqr{2.1} motion & Def.~\ref{def:motion}
 & \eqr{2.34}--\eqr{2.37} & \S\ref{sec:defgrad}, Thm.~\ref{thm:GradgradF}\\
\eqr{2.2} velocity, accel. & Def.~\ref{def:va}
 & \eqr{2.38}--(2.41) & eqs.\ here in \S\ref{sec:defgrad}\\
\eqr{2.3}, \eqr{2.4} reference & Def.~\ref{def:ref}
 & \eqr{2.42}--\eqr{2.44} & Lem.~\ref{lem:dyadprod}\\
\eqr{2.5}--\eqr{2.7} $\bchi_\kappa$ & eqs.\ in \S\ref{sec:pva}
 & \eqr{2.45}, \eqr{2.46} & material curves, \S\ref{sec:defgrad}\\
\eqr{2.8}--\eqr{2.10} velocity forms & \S\ref{sec:pva}
 & \eqr{2.47}--\eqr{2.50} volume & eqs.\ here, Rem.~\ref{rem:fig8}\\
\eqr{2.11}--\eqr{2.13} descriptions & \S\ref{sec:desc}, Rem.~\ref{rem:decorations}
 & \eqr{2.51}--\eqr{2.54} Piola--Nanson & \S\ref{sec:defgrad}\\
\eqr{2.14}--\eqr{2.16} differentials & \S\ref{sec:desc}
 & \eqr{2.55} $\bF^{*}=J\bF^{-t}$ & \S\ref{sec:defgrad}\\
\eqr{2.17}--\eqr{2.21} material deriv. & Thm.~\ref{thm:matder}
 & \eqr{2.56}--\eqr{2.59} mass & \S\ref{sec:mass}\\
\eqr{2.22}--\eqr{2.24} coordinates & \S\ref{sec:cart}, Rem.~\ref{rem:fig6}
 & \eqr{2.60}--\eqr{2.65} localization & Thm.~\ref{thm:localization}\\
\eqr{2.25}--\eqr{2.29} components & \S\ref{sec:cart}
 & Example \eqr{2.30}--\eqr{2.33} & Example~\ref{ex:motion}\\
\eqr{2.66}, \eqr{2.67} stretch & \S\ref{sec:shifter}, Rem.~\ref{rem:fig10}
 & \eqr{2.89}--\eqr{2.95} $\bB$ & Prop.~\ref{prop:B}\\
\eqr{2.68}--\eqr{2.74} shifter & Def.~\ref{def:shifter}, Prop.~\ref{prop:shifter}
 & \eqr{2.96}--\eqr{2.99} strains & Def.~\ref{def:strains}, Rem.~\ref{rem:half}\\
\eqr{2.75}--\eqr{2.78} displacement & \S\ref{sec:shifter}, Rem.~\ref{rem:shifteruse}
 & \eqr{2.100}--\eqr{2.103} arclength & \S\ref{sec:cauchygreen}, Rem.~\ref{rem:fig11}\\
\eqr{2.79}--\eqr{2.81} undeformed & Prop.~\ref{prop:undeformed}
 & \eqr{2.104}--\eqr{2.108} angles & \S\ref{sec:cauchygreen}, Rem.~\ref{rem:fig12}\\
\eqr{2.82}--\eqr{2.88} $\bC$ & Prop.~\ref{prop:C}
 & \S2.8 Problems 1--3 & \S\ref{sec:problems:c2}\\
\hline
\end{tabular}
\end{center}

\medskip\noindent
Figures 2.1--2.12 are described in Remarks~\ref{rem:fig123}, \ref{rem:fig4},
\ref{rem:fig5}, \ref{rem:fig6}, \ref{rem:fig7}, \ref{rem:fig8},
\ref{rem:fig9}, \ref{rem:fig10}, \ref{rem:fig11} and \ref{rem:fig12}. All
algebra uses only Chapter~1: the tensor product and its component rule, the
dyad product $(\ba\otimes\bb)(\bc\otimes\bd)=(\ip\bb\bc)\ba\otimes\bd$, the
transpose rules of Problems 17--20, the box product and determinant
$\tbox{\bF\ba}{\bF\bb}{\bF\bc}=\det\bF\,\tbox\ba\bb\bc$, the cofactor
identity $\bA^{t}\bA^{*}=(\det\bA)\I$ of Problem 22(c), the contracted
$\eps$--$\dd$ identities, and the coordinate free chain rule
$\dl f=\ip{\grad f}{\dl\bx}$.

